\medskip

Soit $a$ un nombre réel strictement supérieur à 1.

On considère la suite $\left(u_n\right)$ définie par $u_0 = a$ et, pour tout entier naturel $n$ : \[u_{n+1} = u_n^2 - 2u_n + 2.\]

On admet que pour tout entier naturel $n$,\: $u_n > 1$.

L'objectif de cet exercice est d'étudier la suite $\left(u_n\right)$ pour différentes valeurs du nombre réel $a$.

\bigskip

\textbf{Partie A : étude de la suite \boldmath $\left(u_n\right)$ \unboldmath dans le cas }\boldmath$1 < a < 2$\unboldmath 

\medskip

\begin{enumerate}
\item 
	\begin{enumerate}
		\item Montrer que, pour tout entier naturel $n$, on a : $u_{n+1} - 2 = u_n\left(u_n - 2\right)$.
		\item Montrer que, pour tout entier naturel $n$, on a : $u_{n+1} - u_n = \left(u_n - 1\right)\left(u_n - 2\right)$.
	\end{enumerate}
\item Dans cette question, on pourra utiliser les égalités établies dans la question précédente.
	\begin{enumerate}
		\item En utilisant un raisonnement par récurrence démontrer que, pour tout entier naturel $n : u_n < 2$.
		\item Montrer que la suite $\left(u_n\right)$ est convergente et déterminer sa limite. 
	\end{enumerate}
\end{enumerate}
	
	
\textbf{Partie B : étude dans le cas particulier }\boldmath $a = 2$\unboldmath

\medskip

\begin{enumerate}
\item ~\\[-25pt]

\begin{minipage}{0.5\linewidth}
On donne ci-contre la fonction \texttt{u} écrite en langage Python.

Déterminer les valeurs renvoyées par le programme lorsque l'on saisit \texttt{u(2,1)} et \texttt{u(2,2)} dans la console Python.
\end{minipage}
\hfill
\begin{minipage}{0.45\linewidth}
{\ttfamily
\begin{tabular}{|l|}\hline
\pmb{def} u(a,n) : \\
\qquad u=a\\
\qquad \pmb{for} k \pmb{in} range(n) :\\
\qquad \quad  u=u**2-2*u+2\\
\qquad \pmb{return} u\\ \hline
\end{tabular}
}
\end{minipage}

\item Quelle conjecture peut-on formuler concernant la suite $\left(u_n\right)$ dans le cas où $a = 2$ ? 

On admettra ce résultat sans démonstration.
\end{enumerate}

\bigskip

\textbf{Partie C : étude dans le cas général}

\medskip

\begin{enumerate}
\item On considère la suite $\left(v_n\right)$ définie, pour tout entier naturel $n$, par $v_n = \ln \left(u_n - 1\right)$.
	\begin{enumerate}
		\item Montrer que la suite $\left(v_n\right)$ est une suite géométrique de raison 2 dont on précisera le premier terme en fonction de $a$.
		\item En déduire que, pour tout entier naturel $n,\: u_n = 1 + \e^{2^n \times \ln (a-1)}$.
	\end{enumerate}
\item Déterminer, suivant les valeurs du réel $a$ strictement supérieur à 1, la limite de la suite $\left(u_n\right)$.
\end{enumerate}

\bigskip
